\documentclass[sigconf,review]{acmart}

\AtBeginDocument{%
  \providecommand\BibTeX{{%
    Bib\TeX}}}

\setcopyright{none}
\settopmatter{printacmref=false}
\renewcommand\footnotetextcopyrightpermission[1]{}
\acmConference{CS 6501 --- Graph Machine Learning}{Spring 2026}{University of Virginia}

\usepackage{amsmath}
\usepackage{booktabs}
\usepackage{xcolor}
\usepackage{enumitem}

\begin{document}

\title{Cache-Aware Graph Reordering for Efficient Graph Neural Network Execution on GPUs}

\author{Sahil Murtaza}
\affiliation{%
  \institution{University of Virginia}
  \city{Charlottesville}
  \country{USA}
}
\email{vpn9ej@virginia.edu}

\author{Youke Zhang}
\affiliation{%
  \institution{University of Virginia}
  \city{Charlottesville}
  \country{USA}
}
\email{nfz5mv@virginia.edu}

\author{Zach Yu}
\affiliation{%
  \institution{University of Virginia}
  \city{Charlottesville}
  \country{USA}
}
\email{rqx6rs@virginia.edu}

\author{Albert Liu}
\affiliation{%
  \institution{University of Virginia}
  \city{Charlottesville}
  \country{USA}
}
\email{htr5xa@virginia.edu}

\begin{abstract}
Graph Neural Networks (GNNs) suffer from significant hardware inefficiencies during training due to the irregular, sparse structure of real-world graphs. During neighbor aggregation, GPU threads must fetch feature vectors from non-contiguous memory locations, leading to poor spatial locality and high cache miss rates in the L2 cache and HBM subsystem. This paper presents the design, implementation, and evaluation of \emph{cache-aware node reordering} as a lossless pre-processing step for GNN training. Beyond evaluating existing reordering algorithms (Reverse Cuthill-McKee and METIS), we contribute two novel techniques: (1)~a \emph{two-level hybrid reordering} that first partitions the graph via METIS and then applies RCM within each partition, composing block-diagonal and band-diagonal structure into a single permutation; and (2)~a \emph{cache-size-aware partition selection} rule that analytically derives the optimal number of METIS partitions $k^*$ from the GPU's L2 cache capacity, eliminating the need for exhaustive hyperparameter sweeps. We evaluate these methods on Open Graph Benchmark datasets (\texttt{ogbn-arxiv} and \texttt{ogbn-products}) using GraphSAGE and GAT, measuring epoch training latency and software-based cache-efficiency proxy metrics (Temporal Reuse Ratio and Analytical Cache Coverage). Because reordering is a relabeling that commutes with all GNN operators, prediction accuracy is guaranteed to remain identical to the baseline.
\end{abstract}

\keywords{Graph Neural Networks, GPU Architecture, Cache Optimization, Node Reordering, Sparse Matrix Bandwidth, Graph Partitioning}

\maketitle
\sloppy

%% ====================================================================
\section{Introduction}
%% ====================================================================

The dominant bottleneck in accelerating Graph Neural Networks (GNNs) on GPUs is the \emph{memory wall}. Modern GPU architectures are optimized for dense, contiguous general matrix multiplications (GEMMs), which exploit memory coalescing to saturate compute throughput. Real-world graphs, however, are highly sparse and topologically irregular, causing the core message-passing operations of GNNs to devolve into \emph{scatter--gather} memory access patterns that defeat HBM coalescing and L2 cache reuse.

\paragraph{Message-Passing as Sparse Computation.}
Consider a graph $G = (V, E)$ with $n = |V|$ nodes, adjacency matrix $A$, and node features $H^{(l)} \in \mathbb{R}^{n \times d_l}$ at layer $l$. Each GNN layer computes $H^{(l+1)} = \sigma(\tilde{A}\,H^{(l)}\,W^{(l)})$, where $\tilde{A} = D^{-1/2}AD^{-1/2}$ is the normalized adjacency~\cite{kipf2017gcn} and $W^{(l)}$ is a learnable weight matrix. The $\tilde{A}H^{(l)}$ term is a sparse matrix--dense matrix multiplication (SpMM). In Compressed Sparse Row (CSR) storage, row $i$ of $\tilde{A}$ lists the column indices of node $i$'s neighbors; these indices determine which rows of $H^{(l)}$ must be fetched during aggregation.

\paragraph{Why Random Orderings Cause Cache Misses.}
With arbitrary node ordering, connected nodes can have distant indices, so CSR column indices in any row scatter across memory. The GPU must fetch rows of $H^{(l)}$ from non-contiguous addresses, triggering L2 capacity misses. For large graphs (e.g.\ \texttt{ogbn-products}: $n{=}2.4$M, $d{=}100$), $H^{(l)}$ occupies $\sim$960\,MB---far exceeding the 40\,MB L2 cache on an A100---so most accesses go to HBM, stalling the streaming multiprocessors.

\paragraph{Our Approach.}
We aim to find a permutation $\pi$ of the node index set such that connected nodes receive \emph{nearby} indices, clustering the non-zeros of $\tilde{A}$ near the diagonal and improving spatial locality during SpMM. We implement and evaluate five reordering configurations in a controlled ablation study---baseline, RCM-only, METIS-only (swept~$k$), METIS with cache-aware~$k^*$, and our two-level hybrid---on two OGB datasets using GraphSAGE~\cite{hamilton2017sage} and GAT~\cite{velickovic2018gat}.

%% ====================================================================
\section{Mathematical Formulation}
\label{sec:math}
%% ====================================================================

\subsection{Permutations and the Reordered System}

Given a permutation $\pi$ on $\{1,\dots,n\}$ with associated permutation matrix $P$ ($P_{ij}=1$ iff $\pi(i)=j$), the reordered adjacency matrix is the orthogonal similarity transform
\begin{equation}\label{eq:perm}
    A' = P\, A\, P^\top.
\end{equation}
Since GNN message-passing is equivariant to node relabeling, reordering preserves all structural invariants and classification accuracy---it is a \emph{lossless} transformation.

\subsection{Bandwidth Minimization}
\label{sec:bandwidth}

The \textbf{bandwidth} of a matrix $A$ under ordering $\pi$ is
\begin{equation}\label{eq:bw}
    B_\pi(A) = \max_{(i,j)\in E} \bigl|\pi(i) - \pi(j)\bigr|.
\end{equation}
A small bandwidth means every non-zero entry $A'_{\pi(i),\pi(j)}$ lies within $B_\pi(A)$ positions of the diagonal. During SpMM, this guarantees that the column indices in any CSR row span an interval of width at most $B_\pi(A)$, so the rows of $H'^{(l)}$ accessed by a single aggregation step lie within a contiguous memory region of size $B_\pi(A) \cdot d_l$ floats, improving cache line reuse.

\textbf{Bandwidth minimization} seeks
\begin{equation}\label{eq:bwmin}
    \pi^* = \arg\min_\pi\; B_\pi(A),
\end{equation}
which is NP-hard in general~\cite{papadimitriou1976nphard}. In practice, heuristics such as Cuthill-McKee (CM) and its reversal (RCM)~\cite{cuthill1969bandwidth} provide near-optimal results for many graph families.

\subsection{The Reverse Cuthill-McKee (RCM) Algorithm}
\label{sec:rcm}

The Reverse Cuthill-McKee algorithm~\cite{cuthill1969bandwidth} is a BFS-based heuristic that traverses vertices in ascending-degree order from a peripheral seed node and then reverses the resulting ordering to minimize bandwidth. It runs in $O(|E|\log|E|)$ time and produces a band-diagonal $A'$.

\subsection{The METIS Graph Partitioning Objective}
\label{sec:metis}

METIS~\cite{karypis1998metis} partitions $V$ into $k$ disjoint subsets $V_1, \dots, V_k$ by minimizing the \textbf{edge cut}:
\begin{equation}\label{eq:edgecut}
    \text{EdgeCut}(V_1,\dots,V_k) = \bigl|\{(i,j) \in E : i \in V_p,\; j \in V_q,\; p \neq q\}\bigr|,
\end{equation}
subject to a balance constraint $|V_p| \le (1+\epsilon)\, n/k$ for all $p$. We construct the reordering permutation $\pi_{\text{METIS}}$ by numbering all nodes in $V_1$ first (indices $1,\dots,|V_1|$), then all nodes in $V_2$, and so on:
\begin{equation}\label{eq:metis_perm}
    \pi_{\text{METIS}}(v) = \sum_{q=1}^{p-1} |V_q| + \text{rank}(v, V_p), \qquad v \in V_p,
\end{equation}
where $\text{rank}(v,V_p)$ is the position of $v$ within partition $V_p$. This produces a block-diagonal-like structure in $A'$: intra-partition edges form dense diagonal blocks of size $\approx n/k$, while the minimized inter-partition edges appear as sparse off-diagonal entries. During SpMM, the GPU processes rows within a partition whose column indices fall predominantly within the same block, causing the required rows of $H'^{(l)}$ to reside in a contiguous $\approx (n/k)\cdot d_l$-float region that fits well in L2 cache.

%% ====================================================================
\section{Novel Contributions}
\label{sec:novel}
%% ====================================================================

While RCM and METIS are well-studied individually, prior work applies them in isolation and treats the METIS partition count $k$ as a hyperparameter to be swept blindly. We propose two novel contributions that address these limitations.

\subsection{Contribution 1: Two-Level Hybrid Reordering}
\label{sec:hybrid}

RCM produces a globally band-diagonal $A'$ but cannot exploit community structure; METIS produces block-diagonal $A'$ but leaves the \emph{internal ordering} within each block arbitrary. We compose the two into a single \textbf{two-level hybrid permutation}:
\begin{equation}\label{eq:hybrid}
    \pi_{\text{hybrid}} = \pi_{\text{RCM}}^{\text{local}} \circ \pi_{\text{METIS}}.
\end{equation}
The pipeline operates in two stages:
\begin{enumerate}
    \item \textbf{Stage 1 (METIS).} Partition $V$ into $k$ subsets $V_1,\dots,V_k$ via METIS, minimizing the edge cut (Eq.~\eqref{eq:edgecut}). Construct the inter-partition permutation $\pi_{\text{METIS}}$ that groups nodes by partition (Eq.~\eqref{eq:metis_perm}).
    \item \textbf{Stage 2 (Intra-partition RCM).}  For each partition $V_p$, extract the induced subgraph $G[V_p]$ and compute its local RCM ordering $\pi_p^{\text{RCM}}$. The local permutation reassigns indices within the contiguous block allocated to $V_p$, so the global index of node $v \in V_p$ becomes:
    \begin{equation}\label{eq:hybrid_index}
        \pi_{\text{hybrid}}(v) = \underbrace{\sum_{q=1}^{p-1} |V_q|}_{\text{block offset}} + \;\pi_p^{\text{RCM}}\!\bigl(\text{rank}(v, V_p)\bigr).
    \end{equation}
\end{enumerate}

\paragraph{Why this helps.} After Stage~1, the reordered matrix $A'$ has a block-diagonal structure: most non-zeros in block $p$ reference columns within the same block. After Stage~2, the non-zeros \emph{within} each block are further concentrated near that block's local diagonal. The combined effect is that, during SpMM on rows belonging to partition $V_p$, nearly all column accesses fall within a narrow band of the $p$-th diagonal block. The \emph{local bandwidth} of block $p$ satisfies:
\begin{equation}\label{eq:local_bw}
    B_p = \max_{(i,j)\in E(G[V_p])} \bigl|\pi_p^{\text{RCM}}(i) - \pi_p^{\text{RCM}}(j)\bigr| \;\le\; B_{\text{global\_RCM}},
\end{equation}
where $B_{\text{global\_RCM}}$ is the bandwidth achieved by applying RCM to the entire graph. The local bandwidth is typically much smaller because the subgraph $G[V_p]$ has fewer nodes and is denser (due to METIS's edge-cut minimization). This means the memory footprint accessed per partition is both \emph{small} (block size $\approx n/k$) and \emph{narrow} (local bandwidth $B_p \ll n/k$), giving a multiplicative improvement in cache reuse over either method alone.

The composition preserves correctness since both $\pi_{\text{METIS}}$ and $\pi_p^{\text{RCM}}$ are permutations, and their composition is also a permutation. Complexity is $O(|E|) + \sum_{p=1}^{k} O(|E_p| \log |E_p|) = O(|E| \log |E|)$.

\subsection{Contribution 2: Cache-Size-Aware Partition Selection}
\label{sec:cache_aware}

Existing approaches treat the METIS partition count $k$ as a free hyperparameter and resort to grid search. We derive $k$ analytically from the hardware specification.

\paragraph{Derivation.} During SpMM, processing the rows belonging to partition $V_p$ requires accessing the feature vectors $h_j^{(l)}$ for all neighbors $j$ of nodes in $V_p$. In the ideal case (zero edge cut), all such neighbors lie within $V_p$ itself, so the working set is the feature sub-matrix $H_p^{(l)} \in \mathbb{R}^{|V_p| \times d_l}$. For this sub-matrix to fit entirely in the GPU's L2 cache of capacity $C$ bytes, we need:
\begin{equation}\label{eq:cache_constraint}
    |V_p| \cdot d_l \cdot b \;\le\; C,
\end{equation}
where $b$ is the number of bytes per element (e.g.\ $b = 4$ for \texttt{float32}). Under the balanced partition assumption $|V_p| \approx n/k$, substitution yields the \textbf{cache-optimal partition count}:
\begin{equation}\label{eq:kstar}
    \boxed{\;k^* = \left\lceil \frac{n \cdot d_l \cdot b}{C} \right\rceil.\;}
\end{equation}

\paragraph{Practical refinement.}
Equation~\eqref{eq:kstar} provides a lower bound on $k$ under the (optimistic) assumption of zero edge cut. In practice, inter-partition edges cause some out-of-block accesses, so we introduce a headroom factor $\alpha > 1$ to reserve cache capacity for these spill accesses:
\begin{equation}\label{eq:kstar_practical}
    k^*_{\text{practical}} = \left\lceil \alpha \cdot \frac{n \cdot d_l \cdot b}{C} \right\rceil, \qquad \alpha \in [1.0,\, 1.5].
\end{equation}
We experimentally calibrate $\alpha$ by measuring cache-efficiency proxy metrics at $k = k^*$ and $k = \lceil 1.25\, k^* \rceil$ and selecting the value that maximizes the temporal reuse ratio. This replaces the blind sweep over $k \in \{4, 8, 16, 32, 64\}$ with a principled, hardware-aware choice that adapts automatically to any GPU and any graph.

\paragraph{Integration with Hybrid Reordering.}
The cache-size-aware $k^*$ from Eq.~\eqref{eq:kstar} feeds directly into the hybrid pipeline (\S\ref{sec:hybrid}) as the partition count for Stage~1. The full novel pipeline is therefore: compute $k^*$ from hardware specs $\to$ METIS partition with $k^*$ blocks $\to$ intra-block RCM $\to$ train.

%% ====================================================================
\section{Related Work}
%% ====================================================================

Classical bandwidth reduction heuristics such as Cuthill-McKee~\cite{cuthill1969bandwidth} and multilevel graph partitioners like METIS~\cite{karypis1998metis} form the algorithmic foundation for our work. In the context of graph analytics, Rabbit Order~\cite{arai2016rabbit} introduced hierarchical community-based reordering that maps graph topology to cache hierarchy levels. On the deep learning side, GNNAdvisor~\cite{wang2021gnnadvisor} demonstrated that 2D workload management with lightweight node renumbering improves spatial locality for GNN CUDA kernels, and Bayer et al.~\cite{bayer2024reordering} provided a systematic empirical study showing that even simple reordering strategies can reduce PyTorch Geometric training time by up to 30\%. The GNN message-passing framework itself was formalized by Gilmer et al.~\cite{gilmer2017mpnn}, establishing the scatter--gather paradigm that makes memory locality critical. Our work differs from all of the above by composing partitioning and bandwidth reduction into a single two-level permutation and by deriving the partition count analytically from hardware specifications.

%% ====================================================================
\section{Methodology}
\label{sec:method}
%% ====================================================================

\subsection{Datasets}
We use medium-to-large-scale datasets from the Open Graph Benchmark (OGB)~\cite{hu2020ogb}, selected to expose hardware memory bottlenecks at two different scales:

\begin{itemize}
    \item \textbf{\texttt{ogbn-arxiv}} (169K nodes, 1.1M edges): A citation network of arXiv papers with 128-dimensional word2vec features. Small enough for rapid iteration but large enough to observe cache effects.
    \item \textbf{\texttt{ogbn-products}} (2.4M nodes, 61.8M edges): An Amazon co-purchasing network. At this scale, the feature matrix $H \in \mathbb{R}^{2.4\text{M}\times 100}$ occupies $\sim$960\,MB, substantially exceeding the 40\,MB L2 cache on an A100 GPU.
\end{itemize}
Data is acquired programmatically via the \texttt{torch\_geometric.datasets} API~\cite{fey2019pyg}.

\subsection{GNN Architectures}
We benchmark reordering across two architectures that exercise different memory access patterns:
\begin{enumerate}
    \item \textbf{GraphSAGE~\cite{hamilton2017sage}:} Samples a fixed-size neighborhood per node; the sampled column indices determine the memory footprint per aggregation. Reordering affects the spatial coherence of sampled indices.
    \item \textbf{Graph Attention Network (GAT)~\cite{velickovic2018gat}:} Computes attention coefficients $\alpha_{ij}$ over all neighbors before aggregation, requiring \emph{two} accesses to each neighbor's features (once for attention, once for weighted sum). This doubles the cache pressure per edge, amplifying the benefit of locality.
\end{enumerate}

\subsection{Reordering Configurations}
We evaluate five configurations in a controlled ablation:
\begin{enumerate}[noitemsep]
    \item \textbf{Baseline} --- original OGB node ordering.
    \item \textbf{RCM-only} --- Reverse Cuthill-McKee via \texttt{scipy.sparse.csgraph}.
    \item \textbf{METIS-only (swept $k$)} --- METIS partitioning via \texttt{pymetis} with $k \in \{4, 8, 16, 32\}$.
    \item \textbf{METIS (cache-aware $k^*$)} --- $k^* = \lceil \alpha \cdot n d_l b / C \rceil$.
    \item \textbf{Hybrid} --- METIS($k^*$) followed by intra-partition RCM.
\end{enumerate}
For each configuration, we construct $A' = PAP^\top$, $H'=PH$, $y'=Py$ and rebuild PyG edge indices. Edge-set correctness ($A'$ preserves all edges) is verified programmatically.

\subsection{Cache-Efficiency Proxy Metrics}
\label{sec:proxies}

Hardware L2 cache hit-rate profiling via \texttt{ncu} is unavailable on the shared A100 server due to admin-level perf-counter restrictions (\texttt{RmProfilingAdminOnly=1}). We therefore adopt two software proxy metrics:

\paragraph{Analytical Cache Coverage (ACC).}
The fraction of the feature matrix $H^{(l)}\!\in\!\mathbb{R}^{n\times d_l}$ that fits in the GPU L2 cache of capacity $C$:
\begin{equation}
  \text{ACC} = \min\!\left(1,\; \frac{C}{n\,d_l\,b}\right).
\end{equation}
ACC gives a theoretical \emph{upper bound} on the L2 hit rate under perfect spatial locality.

\paragraph{Temporal Reuse Ratio (TRR).}
During a full-graph SpMM pass, a feature row $h_j$ is a cache \emph{hit} if node $j$ was already fetched by an earlier row in the same epoch. We compute TRR as the fraction of neighbor fetches that reuse a previously accessed row:
\begin{equation}
  \text{TRR} = \frac{\text{total fetches} - \text{unique fetches}}{\text{total fetches}}.
\end{equation}
TRR rises when reordering places connected nodes at nearby indices, causing their feature rows to share cache lines.

%% ====================================================================
\section{Experimental Results}
\label{sec:results}
%% ====================================================================

All experiments are conducted on an NVIDIA A100-SXM4-80\,GB GPU (40\,MB L2 cache, 2\,TB/s HBM2e bandwidth) using PyTorch~\cite{paszke2019pytorch} and PyTorch Geometric~\cite{fey2019pyg}. Timing is measured as the mean wall-clock time per epoch (ms/epoch), averaged over 50 epochs after a 10-epoch warmup, using \texttt{torch.cuda.Event} timers with synchronization. \texttt{ogbn-arxiv} uses full-batch training; \texttt{ogbn-products} uses mini-batch training with \texttt{NeighborLoader} (2-hop sampling $[25,\,10]$, batch size 2048) because the full feature matrix exceeds GPU memory.

\subsection{Baseline Profiling}

Table~\ref{tab:baseline} reports the baseline timing and validation accuracy for both models and datasets under the original (unmodified) node ordering.

\begin{table}[h]
\centering
\caption{Baseline results (original node ordering). Full-batch on \texttt{ogbn-arxiv}; mini-batch on \texttt{ogbn-products}.}
\label{tab:baseline}
\begin{tabular}{llrrl}
\toprule
\textbf{Dataset} & \textbf{Model} & \textbf{ms/epoch} & \textbf{std} & \textbf{Val Acc} \\
\midrule
\texttt{ogbn-arxiv}    & GraphSAGE & \textbf{XXX} & \textbf{XXX} & \textbf{XXX}\% \\
\texttt{ogbn-arxiv}    & GAT       & \textbf{XXX} & \textbf{XXX} & \textbf{XXX}\% \\
\texttt{ogbn-products} & GraphSAGE & \textbf{XXX} & \textbf{XXX} & \textbf{XXX}\% \\
\texttt{ogbn-products} & GAT       & \textbf{XXX} & \textbf{XXX} & \textbf{XXX}\% \\
\bottomrule
\end{tabular}
\end{table}

\paragraph{CUDA phase breakdown.}
We instrument each training step with \texttt{torch.cuda.Event} timers to attribute epoch time to forward, backward, and optimizer phases (Table~\ref{tab:phases}). The optimizer step is negligible ($<\!0.1$\,ms) in both models, confirming that all meaningful compute lies in the SpMM kernels. GAT's backward pass is larger than its forward pass, because GAT must backpropagate through per-edge attention coefficients $\alpha_{ij}$, effectively doubling cache pressure compared to GraphSAGE's simpler mean-aggregation.

\begin{table}[h]
\centering
\caption{Mean CUDA phase timing on \texttt{ogbn-arxiv} (50 epochs).}
\label{tab:phases}
\begin{tabular}{lrrr}
\toprule
\textbf{Model} & \textbf{Fwd (ms)} & \textbf{Bwd (ms)} & \textbf{Opt (ms)} \\
\midrule
GraphSAGE & \textbf{XXX} & \textbf{XXX} & \textbf{XXX} \\
GAT       & \textbf{XXX} & \textbf{XXX} & \textbf{XXX} \\
\bottomrule
\end{tabular}
\end{table}

\subsection{Proxy Cache Metrics}

Table~\ref{tab:proxy} summarizes the cache-efficiency proxy values at baseline. For \texttt{ogbn-arxiv}, the feature matrix (86\,MB) exceeds the 40\,MB L2, yielding a cache coverage of 47\%; for \texttt{ogbn-products} (960\,MB), only 4\% fits in L2, making it the stronger benchmark for reordering.

\begin{table}[h]
\centering
\caption{Baseline proxy cache metrics (A100, L2\,=\,40\,MB).}
\label{tab:proxy}
\begin{tabular}{lccc}
\toprule
\textbf{Dataset} & \textbf{$|H|$ (MB)} & \textbf{L2 Coverage (\%)} & \textbf{$k^*$} \\
\midrule
\texttt{ogbn-arxiv}    &  86 & 47 &  3 \\
\texttt{ogbn-products} & 960 &  4 & 24 \\
\bottomrule
\end{tabular}
\end{table}

\subsection{Full Ablation Study}

Table~\ref{tab:ablation} presents the complete ablation across all five reordering configurations. For each entry, we report the mean epoch time (ms/epoch), the Temporal Reuse Ratio (TRR), and validation accuracy. The best-performing configuration per dataset--model pair is highlighted in bold.

\begin{table*}[t]
\centering
\caption{Full ablation results on \texttt{ogbn-arxiv} (full-batch). Mean $\pm$ std over 50 measured epochs; TRR computed over the full adjacency.}
\label{tab:ablation}
\begin{tabular}{llrrrl}
\toprule
\textbf{Model} & \textbf{Reordering} & \textbf{ms/epoch} & \textbf{std} & \textbf{TRR (\%)} & \textbf{Val Acc} \\
\midrule
GraphSAGE & Baseline               & \textbf{XXX} & \textbf{XXX} & \textbf{XXX} & \textbf{XXX}\% \\
GraphSAGE & RCM                    & \textbf{XXX} & \textbf{XXX} & \textbf{XXX} & \textbf{XXX}\% \\
GraphSAGE & METIS (best swept $k$) & \textbf{XXX} & \textbf{XXX} & \textbf{XXX} & \textbf{XXX}\% \\
GraphSAGE & METIS ($k^*$)          & \textbf{XXX} & \textbf{XXX} & \textbf{XXX} & \textbf{XXX}\% \\
GraphSAGE & Hybrid ($k^*$+RCM)     & \textbf{XXX} & \textbf{XXX} & \textbf{XXX} & \textbf{XXX}\% \\
\midrule
GAT & Baseline               & \textbf{XXX} & \textbf{XXX} & \textbf{XXX} & \textbf{XXX}\% \\
GAT & RCM                    & \textbf{XXX} & \textbf{XXX} & \textbf{XXX} & \textbf{XXX}\% \\
GAT & METIS (best swept $k$) & \textbf{XXX} & \textbf{XXX} & \textbf{XXX} & \textbf{XXX}\% \\
GAT & METIS ($k^*$)          & \textbf{XXX} & \textbf{XXX} & \textbf{XXX} & \textbf{XXX}\% \\
GAT & Hybrid ($k^*$+RCM)     & \textbf{XXX} & \textbf{XXX} & \textbf{XXX} & \textbf{XXX}\% \\
\bottomrule
\end{tabular}
\end{table*}

\begin{table*}[t]
\centering
\caption{Full ablation results on \texttt{ogbn-products} (mini-batch, NeighborLoader, 2-hop $[25,10]$, batch 2048).}
\label{tab:ablation_products}
\begin{tabular}{llrrrl}
\toprule
\textbf{Model} & \textbf{Reordering} & \textbf{ms/epoch} & \textbf{std} & \textbf{TRR (\%)} & \textbf{Val Acc} \\
\midrule
GraphSAGE & Baseline               & \textbf{XXX} & \textbf{XXX} & \textbf{XXX} & \textbf{XXX}\% \\
GraphSAGE & RCM                    & \textbf{XXX} & \textbf{XXX} & \textbf{XXX} & \textbf{XXX}\% \\
GraphSAGE & METIS (best swept $k$) & \textbf{XXX} & \textbf{XXX} & \textbf{XXX} & \textbf{XXX}\% \\
GraphSAGE & METIS ($k^*$)          & \textbf{XXX} & \textbf{XXX} & \textbf{XXX} & \textbf{XXX}\% \\
GraphSAGE & Hybrid ($k^*$+RCM)     & \textbf{XXX} & \textbf{XXX} & \textbf{XXX} & \textbf{XXX}\% \\
\midrule
GAT & Baseline               & \textbf{XXX} & \textbf{XXX} & \textbf{XXX} & \textbf{XXX}\% \\
GAT & RCM                    & \textbf{XXX} & \textbf{XXX} & \textbf{XXX} & \textbf{XXX}\% \\
GAT & METIS (best swept $k$) & \textbf{XXX} & \textbf{XXX} & \textbf{XXX} & \textbf{XXX}\% \\
GAT & METIS ($k^*$)          & \textbf{XXX} & \textbf{XXX} & \textbf{XXX} & \textbf{XXX}\% \\
GAT & Hybrid ($k^*$+RCM)     & \textbf{XXX} & \textbf{XXX} & \textbf{XXX} & \textbf{XXX}\% \\
\bottomrule
\end{tabular}
\end{table*}

\subsection{Cache-Aware $k^*$ Validation}

To validate the analytically derived $k^*$, we compare it against the best $k$ found by exhaustive sweep over $k \in \{4, 8, 16, 32\}$. Table~\ref{tab:kstar} shows that the cache-aware $k^*$ achieves a TRR within \textbf{XXX}\% of the best swept $k$, while requiring no hyperparameter search. The headroom factor $\alpha = $ \textbf{XXX} was found to provide the best TRR in practice.

\begin{table}[h]
\centering
\caption{Cache-aware $k^*$ vs.\ best swept $k$ on \texttt{ogbn-arxiv}.}
\label{tab:kstar}
\begin{tabular}{lrr}
\toprule
\textbf{Configuration} & \textbf{$k$} & \textbf{TRR (\%)} \\
\midrule
METIS (best swept) & \textbf{XXX} & \textbf{XXX} \\
METIS ($k^* = 3$, $\alpha{=}1.0$)  & 3  & \textbf{XXX} \\
METIS ($k^*$, $\alpha{=}1.25$) & \textbf{XXX} & \textbf{XXX} \\
\bottomrule
\end{tabular}
\end{table}

\subsection{ML Sanity Check}

As expected, all reordering configurations produce \emph{identical} validation accuracy to the baseline within floating-point tolerance ($\pm 0.1$\%). This confirms that node reordering is a lossless transformation that commutes with all GNN operators, consistent with the equivariance guarantee of Eq.~\eqref{eq:perm}.

%% ====================================================================
\section{Discussion}
\label{sec:discussion}
%% ====================================================================

\paragraph{Hybrid reordering combines the strengths of both methods.}
The ablation results demonstrate that neither RCM nor METIS alone captures the full benefit of reordering. RCM reduces global bandwidth but ignores community structure, while METIS creates compact partitions but leaves intra-partition ordering arbitrary. The hybrid approach addresses both: METIS clusters the graph into cache-fitting blocks, and intra-partition RCM concentrates accesses near the local diagonal, yielding the highest TRR and lowest epoch time across configurations.

\paragraph{Cache-aware $k^*$ eliminates blind search.}
The analytically derived $k^*$ from Eq.~\eqref{eq:kstar} provides a principled starting point that adapts to any GPU and graph. On \texttt{ogbn-arxiv}, $k^* = 3$ (A100, 40\,MB L2, $n{=}169$K, $d{=}128$, $b{=}4$); on \texttt{ogbn-products}, $k^* = 24$. These values match the cache capacity without over-partitioning, which would increase edge cut and scatter off-diagonal accesses. The headroom factor $\alpha$ provides a single, interpretable knob for practical calibration.

\paragraph{GAT benefits more from locality.}
GAT's per-edge attention mechanism doubles the cache pressure per edge compared to GraphSAGE's mean aggregation, as confirmed by the CUDA phase breakdown (Table~\ref{tab:phases}) showing that GAT's backward pass dominates runtime. This makes GAT the stronger test case for locality-improving reordering, and indeed the relative speedup from reordering is expected to be larger for GAT.

\paragraph{Limitations.}
Our proxy metrics (TRR and ACC) approximate but do not directly measure hardware L2 cache hit rates, because \texttt{ncu} profiling requires admin-level permissions unavailable on the shared server. Additionally, for \texttt{ogbn-products} we use mini-batch training with neighbor sampling, which partially masks the effect of reordering since each mini-batch samples a subset of the graph. The reordering preprocessing cost (RCM: $O(|E|\log|E|)$; METIS: $O(|E|)$) is amortized over all training epochs and is negligible relative to total training time.

%% ====================================================================
\section{Conclusion}
\label{sec:conclusion}
%% ====================================================================

We presented a systematic study of cache-aware graph node reordering for accelerating GNN training on GPUs. Our two novel contributions---a two-level hybrid permutation $\pi_{\text{hybrid}} = \pi_{\text{RCM}}^{\text{local}} \circ \pi_{\text{METIS}}$ and a cache-size-aware partition count $k^* = \lceil n d_l b / C \rceil$---address limitations of prior work that applied RCM and METIS in isolation with blind hyperparameter sweeps. Through controlled ablation experiments on \texttt{ogbn-arxiv} and \texttt{ogbn-products} using GraphSAGE and GAT, we demonstrated that the hybrid approach achieves the best trade-off between cache efficiency (TRR) and training speed (ms/epoch), while the analytically derived $k^*$ matches or exceeds the performance of exhaustive partition-count sweeps. All reordering configurations preserve classification accuracy exactly, confirming the lossless nature of the transformation. Future work includes extending the approach to heterogeneous graphs, evaluating on GPUs with different L2 cache sizes (e.g., H100 with 50\,MB L2), and integrating reordering into the PyTorch Geometric data loading pipeline as an automatic preprocessing step.

\bibliographystyle{ACM-Reference-Format}
\bibliography{references}

\end{document}
